\chapter{Kan extensions}
\label{chap:kanextensions}

The purpose of this chapter is to introduce the language of 2-categories and the apparatus of Kan extensions.

\section{The language of 2-categories}
\label{sec:functorsnaturammaps}

\index{2-category}
Recall that a 2-category ${\cat C}$ is a category enriched over categories. This means by definition that for each two objects $A, B$ of ${\cat C}$ the maps $A \ra B$ form the objects of a category ${\cat C}(A, B)$. The the composition functor $c_{ABC} \colon {\cat C}(A, B) \times {\cat C}(B, C) \ra {\cat C}(A, C)$ is required to be associative and to have $1_A \colon * \ra {\cat C}(A, A)$ as a left and right unit, where $*$ denotes the point-category. 

The objects of a 2-category ${\cat C}$ are called 0-cells, the objects of ${\cat C}(A, B)$ are called 1-cells and the morphisms of ${\cat C}(A, B)$ are called 2-cells. A good introduction to 2-categories can be found in Kelly-Street \cite{Kelly-Street1} or Borceux \cite{Borceux}.

This section describes the notation we use for compositions of 1-cells and 2-cells in a 2-category. Each notation has a full form and a simple form. The simple form of the notation is ambiguous, and is only used if it is clear from the context which functor or natural map operation we refer to. 

We denote as usual 1-cells $f: A \ra B$ with a single arrow. Between 1-cells $f, g \colon A \ra B$, we denote 2-cells as $\alpha \colon f \Ra g$, or just $\alpha \colon f \ra g$ if no confusion can occur. 

The composition of 1-cells $f, g$ 
\begin{equation}
\label{eqn:comp1}
\xymatrix{
    A \ar[r]^f &
    B \ar[r]^g &
    C
}
\end{equation}
is just the composition at the level of unenriched hom-sets ${\cat C}(-, -)$ and is denoted $g \comp f \colon A \ra C$, or in simple form $gf$.

The composition of 2-cells $\alpha, \beta$ 
\begin{equation}
\label{eqn:comp2}
\xymatrix{
    A \ar@<4ex>[rr]^(.3)f _\,="f" \ar[rr]^(.3)g ^\,="g1" _\,="g2"  \ar@<-4ex>[rr] ^(.3)h^\,="h" & &
    B
    \ar@{=>} "f";"g1" ^{\alpha}
    \ar@{=>} "g2";"h" ^{\beta}
}
\end{equation}
is composition at the level of ${\cat C}(-, -)$ and is denoted $\beta \odot \alpha \colon f \Ra h$, or in simple form $\beta \alpha$.

The composition of a 2-cell $\alpha$ with a 1-cell $f$
\begin{equation}
\label{eqn:comp3}
\xymatrix{
    A \ar[rr]^f & &
    B \ar@<2ex>[rr]^(.3)g _\,="g" \ar@<-2ex>[rr] ^(.3)h^\,="h" & & 
    C
    \ar@{=>} "g";"h" ^{\alpha}
}
\end{equation}
is just $c_{ABC}(1_f, \alpha)$, and we denote it $\alpha . f \colon gf \Ra hf$ or in simple form $\alpha f$. The composition in the other direction
\begin{equation}
\label{eqn:comp4}
\xymatrix{
    A \ar@<2ex>[rr]^(.3)f _\,="f" \ar@<-2ex>[rr] ^(.3)g^\,="g" & &
    B \ar[rr]^h & & 
    C
    \ar@{=>} "f";"g" ^{\alpha}
  }
\end{equation}
is $c_{ABC}(\alpha, 1_h)$ and we denote it $h . \alpha \colon hf \Ra hg$, or in simple form $h \alpha$.

The notations we have established up until now can be used to completely describe compositions of 1- and 2-cells. However it is convenient to introduce the $\star$ notation to denote the composition of adjacent of 2-cells of planar diagrams.

We will denote the composition of 2-cells $\alpha \colon jf \Ra g$ and $\beta \colon h \Ra ij$
\begin{equation}
\label{eqn:comp5}
\xymatrix {
    & B \ar[dr]^h \ar[dd]^j & \\
    A \ar[ur]^f \ar[dr]_g \ar@{}[rr]|(.3){\big\Downarrow \alpha}|(.7){\big\Downarrow \beta} & & D \\
    & C \ar[ur]_i &
  }
\end{equation}
as $\beta \star \alpha = (i . \alpha) \odot (\beta . f)$.

We use the same $\star$ notation to denote the composition of 2-cells $\alpha \colon f \Ra jg$ and $\beta \colon hj \Ra i$
\begin{equation}
\label{eqn:comp6}
\xymatrix {
    & B \ar[dr]^h & \\
    A \ar[ur]^f \ar[dr]_g \ar@{}[rr]|(.3){\big\Downarrow \alpha}|(.7){\big\Downarrow \beta} & & D  \\
    & C \ar[ur]_i \ar[uu]_j &
  }
\end{equation}
as $\beta \star \alpha = (\beta . g) \odot (h . \alpha)$.

In particular, taking $j$ to be an identity map in (\ref{eqn:comp5}) or (\ref{eqn:comp6}) we get the composition of 2-cells $\beta \star \alpha = c_{ABC}(\alpha, \beta)$
\begin{equation}
\label{eqn:comp7}
\xymatrix{
    A \ar@<2ex>[rr]^(.3)f _\,="f" \ar@<-2ex>[rr] ^(.3)g^\,="g" & &
    B \ar@<2ex>[rr]^(.3)h _\,="h" \ar@<-2ex>[rr] ^(.3)i^\,="i" & & 
    C
    \ar@{=>} "f";"g" ^{\alpha}
    \ar@{=>} "h";"i" ^{\beta}
  }
\end{equation}
Since $c_{ABC}$ is a functor, $(i . \alpha) \odot (\beta . f) = (\beta . g) \odot (h . \alpha)$ and the $\star$ notation is consistent in (\ref{eqn:comp7}) with (\ref{eqn:comp5}) and (\ref{eqn:comp6}). 

In simple form, if no confusion is possible we denote $\beta \alpha$ for $\beta \star \alpha$.

There is a general theorem called the Pasting Theorem regarding compositions of 2-cells of planar diagrams, for which we refer to \cite{Power} . 

\index{2-category!opposites ${\cat C}^{1-op}$, ${\cat C}^{2-op}$, ${\cat C}^{1,2-op}$}
A 2-category ${\cat C}$ admits three flavours of opposites, denoted respectively:
\begin{enumerate}
\item
${\cat C}^{1-op}$ (also denoted ${\cat C}^{co}$), reverting the direction of 1-cells
\item
${\cat C}^{2-op}$ (also denoted ${\cat C}^{op}$), reverting the direction of 2-cells
\item
${\cat C}^{1,2-op}$ (also denoted ${\cat C}^{coop}$), reverting the direction of both 1- and 2-cells
\end{enumerate}

\index{2-category!2-functor}
\index{2-category!2-functor!lax or right weak}
\index{2-category!2-functor!op-lax or left weak}
\index{2-category!2-functor!pseudo 2-functor}
\index{2-category!2-functor!strict}
Let us also recall the various flavours of 2-functors $F \colon {\cat C}_1 \ra {\cat C}_2$. All flavours send the 0, 1 and 2-cells of ${\cat C}_1$ respectively to 0, 1 and 2-cells of ${\cat C}_2$, and preserve the composition and units of 2-cells on the nose. The functor $F$ is called

\begin{enumerate}
\item
{\it lax} or {\it right weak} if it preserves the composition and units of 1-cells up to canonical 2-cells $\Gamma_{f, g} \colon  F_gF_f \Ra F_{gf}$ for $A \overset{f}{\ra} B \overset{g}{\ra} C$ in ${\cat C}_1$ and $\Delta_A \colon 1_{F_A} \Ra F_{1_A}$ for $A$ an object of ${\cat C}_1$, so that if $c \overset{h}{\ra} d$ in ${\cat C}_1$ then the diagrams below are commutative
\[
\xymatrix{
  F_hF_gF_f \ar[d]_{F_h \star \Gamma_{f, g}} \ar[rr]^{\Gamma_{g, h} \star F_f} && F_{hg}F_f \ar[d]^{\Gamma_{f, hg}} \\
  F_hF_{gf} \ar[rr]_{\Gamma_{gf, h}} && F_{hgf}
}
\]
\[
\xymatrix{
  F_f \comp 1_{F_A} \ar[dr]_{=} \ar[r]^{F_f \star \Delta_A} & F_f \comp F_{1_A} \ar[d]^{\Gamma_{1_A, f}} && F_{1_B} \comp F_{f} \ar[d]_{\Gamma_{f, 1_B}} & 1_{F_B} \comp F_f \ar[l]_{\Delta_B \star F_f} \ar[ld]^{=} \\
  & F_{f \comp 1_A} && F_{1_B \comp f} &
}
\]
$\Gamma_{f, g}$ are called {\it composition} 2-cells and $\Delta_A$ are called {\it unit} 2-cells.
\item
{\it op-lax} or {\it left weak} if composition and unit of 1-cells are functorial up to natural 2-cells $\Gamma^{'}_{f, g}, \Delta^{'}_A$ going in the opposite direction, making commutative three diagrams dual to the ones listed above
\item
a {\it pseudo} functor if it is lax and all $\Gamma_{f, g}$, $\Delta_A$ are {\it isomorphisms}. Notice that this is equivalent to saying that $F$ is op-lax and  all $\Gamma^{'}_{f, g}$, $\Delta^{'}_A$ are isomorphisms. A pseudo-functor thus preserves composition and units of 1-cells up to canonical isomorphisms $\Gamma_{f, g}$, $\Delta_A$.
\item
{\it strict} if it is a pseudo-functor whose all $\Gamma_{f, g}$, $\Delta_A$ are {\it identities}. A strict 2-functor thus preserves composition and units of 1-cells on the nose.
\end{enumerate}

If $F \colon {\cat C}_1 \ra {\cat C}_2$ is lax or op-lax, we will also say that $F$ is
\begin{enumerate}
\item
{\it pseudo-unital} if all its unit 2-cells $\Delta_a \colon 1_{Fa} \Ra F_{1a}$ (or $\Delta^{'}_a$) are {\it isomorphisms}
\item
{\it strictly unital}, if all its unit 2-cells $\Delta_a \colon 1_{Fa} \Ra F_{1a}$ (or $\Delta^{'}_a$) are {\it identities}
\end{enumerate} 

A pseudo-functor is pseudo-unital, and a strict 2-functor is strictly unital\footnote{In fact, any lax pseudo-unital (or op-lax pseudo-unital, or pseudo-) functor can be ``rectified'' to a strictly unital lax (resp. strictly unital op-lax, strictly unital pseudo-) functor. We are purposefully vague on this point, to avoid having to talk about natural transformations of 2-functors.}.

A 2-subcategory ${\cat C}^{'}$ of a 2-category ${\cat C}$ consists of a subclass of objects $Ob {\cat C}^{'} \subset Ob{\cat C}$ along with subcategories ${\cat C}^{'}(A, B) \subset {\cat C}(A, B)$ for objects $A, B \eps Ob {\cat C}^{'}$ that are stable under the composition rule $c_{ABC}$ for $A, B, C \eps Ob{\cat C}^{'}$ and include the image of the unit $1_A \colon * \ra {\cat C}(A, A)$ for $A \eps Ob{\cat C}^{'}$. A 2-subcategory is 2-{\it full} if ${\cat C}^{'}(A, B) = {\cat C}(A, B)$ for any $A, B \eps Ob {\cat C}^{'}$.

The category of categories $CAT$ forms a 2-category, denoted $2CAT$, with categories as 0-cells, functors as 1-cells and natural maps as 2-cells. We will use the notation introduced in this section to denote compositions of functors and natural maps. In particular, natural maps will be denoted with a double arrow $\Ra$ (or simply with $\ra$ if no confusion is possible), and composition of natural maps viewed as adjacent 2-cells will be denoted with $\star$. 

The category of small categories forms a 2-subcategory of $2CAT$, denoted $2{\cat Cat}$.

Any category ${\cat C}$ is in a canonical way a 2-category, having only identity 2-cells. In particular, we can talk of lax, op-lax, pseudo and strict 2-functors ${\cat C} \ra 2CAT$. A strict 2-functor ${\cat C} \ra 2CAT$ is the same as a functor ${\cat C} \ra CAT$.

For the rest of the chapter, we limit ourselves to the case of $2CAT$. As an exercise, the reader may enjoy reformulating the definitions and the results that follow so that they make sense in an arbitrary 2-category (cf. \cite{Kelly-Street1}, \cite{Street1}).

\section{Adjoint functors}
\label{sec:adjointfunctors}
\index{functor!adjoint}
Recall that an adjunction $u_1 \dashv u_2$ between two functors $u_1 : {\cat A} \rightleftarrows {\cat B} : u_2$ is a bijection of sets
\begin{center}
$\zeta \colon {\cat B}(u_1A, B) \cong {\cat A}(A, u_2B)$ 
\end{center}
natural in objects $A \eps {\cat A}$ and $B \eps {\cat B}$. For example, if $u_1$, $u_2$ are equivalences of categories, then $u_1 \dashv u_2$ (and $u_2 \dashv u_1$).

Whenever we say that $u_1, u_2$ is an adjoint pair we refer to a particular bijection $\zeta$. The following proposition encodes the basic properties of adjoint functors that we will need.

\begin{prop}
\label{prop:adjfunctors}
Suppose that $u_1 : {\cat A} \rightleftarrows {\cat B} : u_2$ is a pair of functors.
\begin{enumerate}
\item
An adjunction $u_1 \dashv u_2$ is uniquely determined by natural maps $\phi \colon 1_{\cat A} \Ra u_2u_1$ (the {\it unit}) and $\psi \colon u_1u_2 \Ra 1_{\cat B}$ (the {\it counit} of the adjunction) with the property that both the following compositions are identities
\begin{center}
$\xymatrix {
    {u_1} \ar@{=>}[r]^-{{u_1} {\phi}} & {u_1}{u_2}{u_1} \ar@{=>}[r]^-{{\psi} {u_1}} & {u_1} & {u_2} \ar@{=>}[r]^-{{\phi} {u_2}} & {u_2}{u_1}{u_2} \ar@{=>}[r]^-{{u_2} {\psi}} & {u_2} 
  }
$
\end{center}
\item
If $u_1 \dashv u_2$ is an adjunction with unit $\phi$ and counit $\psi$, then
\begin{enumerate}
\item
$u_1$ (resp. $u_2$) is fully faithful iff $\phi$ (resp. $\psi$) is a natural isomorphism.
\item
$u_1$ and $u_2$ are inverse equivalences of categories iff both $\phi$ and $\psi$ are natural isomorphisms.
\end{enumerate}
\end{enumerate}
\end{prop}

\begin{proof}
See for example Mac Lane \cite{MacLane}.
\end{proof}

Here is another way to state part (2) of the previous proposition. The proof is left to the reader.
\begin{prop}
\label{prop:adjfunctors2}
Suppose that $u_1 \dashv u_2$ is an adjoint pair of functors as above. Then the following statements are equivalent:
\begin{enumerate}
\item[(1)]
(resp. (1r), resp. (1l)). 
For any objects $A \eps {\cat A}$, $B \eps {\cat B}$, a map $u_1A \ra B$ is an isomorphism iff (resp. if, resp. only if) its adjoint $A \ra u_2B$ is an isomorphism
\item[(2)]
(resp. (2r), resp. (2l)). 
$u_1$ and $u_2$ are inverse equivalences of categories (resp. $u_2$ is fully faithful, resp $u_1$ is fully faithful).
\end{enumerate}
\end{prop}


\section{Kan extensions}
\label{sec:leftrightkanextension}

\index{functor!derived}
\index{functor!kan extension}
\begin{defn}
\label{defn::def:kanext}
Consider two functors $u \colon {\cat A} \ra {\cat B}$ and $\gamma \colon {\cat A} \ra {\cat A}^{'}$
\begin{enumerate}
\item
A left Kan extension of $u$ along $\gamma$ is a pair $({\bf L}_\gamma u, \eps)$ where ${\bf L}_\gamma u \colon {\cat A}^{'} \ra {\cat B}$ is a functor and $\eps \colon {\bf L}_\gamma u \comp \gamma \Ra u$ is a natural map
\begin{center}
$\xymatrix {
    {\cat A} \ar[rrd]^u \ar[d]_{\gamma} &  \\ 
    {\cat A}^{'} \ar[rr]_{{\bf L}_\gamma u} \ar@{}[ru]|(.42){\;\; \eps \; \Uparrow } && {\cat B}
  }$
\end{center}
satisfying the following universal property: if $(\Lambda, \lambda)$ is another pair of a functor $\Lambda \colon {\cat A}^{'} \ra {\cat B}$ and natural map $\lambda \colon \Lambda \comp \gamma \Ra u$
\begin{center}
$\xymatrix {
    {\cat A} \ar@<1ex>[rrd]^u \ar[d]_{\gamma} &&  \\ 
    {\cat A}^{'} \ar@/^1pc/[rr]_(.6){{\bf L}_\gamma u}_(.4)\,="d" \ar@/_1pc/[rr]_(.6){\Lambda} ^(.4)\,="c" \ar@{}[ru]|(.45){\;\; \eps \; \Uparrow } && {\cat B}
    \ar@{} "c";"d" ^{\delta \; \Uparrow}
  }$
\end{center}
then there exists a unique natural map $\delta \colon \Lambda \Ra {\bf L}_\gamma u$ with $\eps \star \delta = \lambda$
\item
A right Kan extension of $u$ along $\gamma$ is a pair $({\bf R}_\gamma u, \nu)$ where ${\bf R}_\gamma u \colon {\cat A}^{'} \ra {\cat B}$ is a functor and $\nu \colon u \Ra {\bf R}_\gamma u \comp \gamma$ is a natural map
\begin{center}
$\xymatrix {
    {\cat A} \ar[rrd]^u \ar[d]_{\gamma} &&  \\ 
    {\cat A}^{'} \ar[rr]_{{\bf R}_\gamma u} \ar@{}[ru]|(.42){\;\; \nu \; \Downarrow } & & {\cat B}
  }$
\end{center}
satisfying the following universal property: if $(\Lambda, \lambda)$ is another pair of a functor $\Lambda \colon {\cat A}^{'} \ra {\cat B}$ and a natural map $\lambda \colon u \Ra \Lambda \comp \gamma $
\begin{center}
$\xymatrix {
    {\cat A} \ar@<1ex>[rrd]^u \ar[d]_{\gamma} &&  \\ 
    {\cat A}^{'} \ar@/^1pc/[rr]_(.6){{\bf R}_\gamma u}_(.4)\,="d" \ar@/_1pc/[rr]_(.6){\Lambda} ^(.4)\,="c" \ar@{}[ru]|(.45){\;\; \nu \; \Downarrow } && {\cat B}
    \ar@{} "c";"d" ^{\delta \; \Downarrow}
  }$
\end{center}
then there exists a unique natural map $\delta \colon {\bf R}_\gamma u \Ra \Lambda$ with $\delta \star \nu = \lambda$
\end{enumerate}
\end{defn}

The left Kan extension $({\bf L}_\gamma u, \eps)$ is also called in some references the {\it left derived} of $u$ along $\gamma$ and the right Kan extension $({\bf R}_\gamma u, \eps)$ the {\it right derived} of $u$ along $\gamma$. Since it is defined by an universal property, if the left (or right) Kan extension exists then it is unique up to a unique isomorphism. 

The next proposition is an existence criterion for Kan extensions: if $\gamma$ admits a left (resp. right) adjoint, then $u$ admits a left (resp. right) Kan extension along $\gamma$.

\begin{prop}
\label{prop:existencekanextension}
Consider two functors $u \colon {\cat A} \ra {\cat B}$ and $\gamma \colon {\cat A} \ra {\cat A}^{'}$.
\begin{enumerate}
\item
If $\gamma$ admits a left adjoint $\gamma^{'}$ with adjunction unit $\phi \colon 1_{{\cat A}^{'}} \Ra \gamma\gamma^{'}$ and counit $\psi \colon \gamma^{'}\gamma \Ra 1_{{\cat A}}$, then $(u\gamma^{'}, u \psi)$ is a left Kan extension of $u$ along $\gamma$
\item
If $\gamma$ admits a right adjoint $\gamma^{'}$ with adjunction unit $\phi \colon 1_{\cat A} \Ra \gamma^{'}\gamma$ and counit $\psi \colon \gamma\gamma^{'} \Ra 1_{{\cat A}^{'}}$, then $(u\gamma^{'}, u \phi)$ is a right Kan extension of $u$ along $\gamma$
\end{enumerate}
\end{prop}

\begin{proof}
We only prove (1). 
\begin{center}
$\xymatrix {
    {\cat A} \ar@<1ex>[rrd]^u \ar[d]_{\gamma} &&  \\ 
    {\cat A}^{'} \ar@/^1pc/[rr]_(.6){u \gamma^{'}}_(.4)\,="d" \ar@/_1pc/[rr]_(.6){\Lambda} ^(.4)\,="c" \ar@{}[ru]|(.45){\;\; u\psi \; \Uparrow } \ar@<2ex>[u]^(.55){\gamma^{'}} && {\cat B}
    \ar@{} "c";"d" ^{\delta \; \Uparrow}
  }$
\end{center}

For any pair $(\Lambda, \lambda)$ with $\Lambda \colon {\cat A}^{'} \ra {\cat B}$ and $\lambda \colon \Lambda \gamma \Ra u$, we'd like to show that there exists a unique natural map $\delta \colon \Lambda \Ra u \gamma^{'}$ with
\begin{equation}
\label{eqn:existencekanextension1}
\lambda = (u \psi) \comp (\delta \gamma)
\end{equation}
To show existence, we define $\delta \colon \Lambda \Ra \Lambda \gamma \gamma^{'} \Ra u \gamma^{'}$ as the composition 
\begin{equation}
\label{eqn:existencekanextension2}
\delta = (\lambda \gamma^{'}) \comp (\Lambda \phi)
\end{equation}
$\delta$ defined by (\ref{eqn:existencekanextension2}) satisfies (\ref{eqn:existencekanextension1}), using the commutative diagram
\begin{center}
$\xymatrix{
    \Lambda \gamma \ar@{=>}[r]^-{\Lambda \phi \gamma} \ar@{=>}[rd]_{1_{\Lambda \gamma}} & \Lambda \gamma \gamma^{'} \gamma \ar@{=>}[r]^{\lambda \gamma^{'} \gamma} \ar@{=>}[d]_{\Lambda \gamma \psi} & u \gamma^{'} \gamma \ar@{=>}[d]^{u\psi} \\
    & \Lambda \gamma \ar@{=>}[r]^{\lambda} & u
    }$
\end{center} 
To show uniqueness, if a map $\delta$ satisfies (\ref{eqn:existencekanextension1}) then $\delta$ satisfies (\ref{eqn:existencekanextension2}) because of the commutative diagram
\begin{center}
$\xymatrix{
    \Lambda  \ar@{=>}[d]_-{\Lambda \phi} \ar@{=>}[r]^{\delta} & u \gamma^{'} \ar@{=>}[d]_{u \gamma^{'} \phi} \ar@{=>}[rd]^{1_{u \gamma^{'}}} &  \\
     \Lambda \gamma \gamma^{'} \ar@{=>}[r]_{\delta \gamma \gamma^{'}} & u \gamma^{'} \gamma \gamma^{'} \ar@{=>}[r]_{ u \psi \gamma^{'}} & u \gamma^{'}
    }$
\end{center} 
\end{proof}

Kan extensions commute with composition {\it along the base functor} $\gamma$, as seen in the next propoition.

\begin{prop}
\label{prop:basecompositionkanextension}
Consider the functors ${\cat A} \overset{u}{\ra} {\cat B}$, ${\cat A} \overset{\gamma}{\ra} {\cat A}^{'} \overset{\gamma^{'}}{\ra} {\cat A}^{''}$.
\begin{enumerate}
\item
Suppose that $({\bf L}_\gamma u, \eps)$ exists. Then $({\bf L}_{\gamma^{'}} {\bf L}_\gamma u, \eps^{'})$ exists iff $({\bf L}_{\gamma^{'}\gamma} u, \eps^{''})$ exists. If they both exist, then the latter is isomorphic to $({\bf L}_{\gamma^{'}}{\bf L}_\gamma u, \eps \star \eps^{'})$.
\item
Suppose that $({\bf R}_\gamma u, \nu)$ exists. Then $({\bf R}_{\gamma^{'}} {\bf R}_\gamma u, \nu^{'})$ exists iff $({\bf R}_{\gamma^{'}\gamma} u, \nu^{''})$ exists, in which case the latter is isomorphic to $({\bf R}_{\gamma^{'}}{\bf R}_\gamma u, \nu^{'} \star \nu)$.
\end{enumerate}
\end{prop}

\begin{proof}
Immediate using the universal property of the left (resp. right) Kan extensions. 
%\[
%\xymatrix{
%  {\cat A} \ar[rr]^u \ar[d]_{\gamma} && {\cat B} \\
%  {\cat A}^{'} \ar[d]_{\gamma^{'}} && \\
%  {\cat A}^{''} &&
%}
%\]
\end{proof}

We state a corollary needed in the proof of \theoremref{thm:genholim}.
\begin{cor}
\label{cor:neededbygenhomlim}
Consider two functors $u \colon {\cat A} \ra {\cat B}$ and $\gamma \colon {\cat A} \ra {\cat A}^{'}$, and a pair of inverse equivalences of categories $\gamma^{'} : {\cat A}^{'} \leftrightarrows {\cat A}^{''} : \gamma^{''}$, with natural isomorphisms $\phi \colon 1_{{\cat A}^{'}} \Ra \gamma^{''}\gamma^{'}$ and $\psi \colon 1_{{\cat A}^{''}} \Ra \gamma^{'}\gamma^{''}$.
\begin{enumerate}
\item
The left Kan extension $({\bf L}_\gamma u, \eps)$ exists iff the left Kan extension  $({\bf L}_{{\gamma}^{'}\gamma} u, \eps^{'})$ exists. If they both exist, then the latter is isomorphic to $(({\bf L}_\gamma u) \gamma^{''}$, $\eps \star \psi)$.
\item
The right Kan extension $({\bf R}_\gamma u, \nu)$ exists iff the right Kan extension $({\bf R}_{{\gamma}^{'}\gamma} u$, $\nu^{'})$ exists, in which case the latter is isomorphic to $(({\bf R}_\gamma u) \gamma^{''}$, $\phi \star \nu)$.
\end{enumerate}
\end{cor}

\begin{proof}
Consequence of \propositionref{prop:existencekanextension} and \propositionref{prop:basecompositionkanextension}.
\end{proof}

